\clearpage
\fancyhead[L]{\small BODE. ASTRONOMISCHES JAHRBUCH.}
\hypertarget{native-581-lower}{}
\begin{center}\small Göttingische gelehrte Anzeigen. Stück 28. Seite 265..277. 1815 Februar 18.\\[3pt]\rule{.72\linewidth}{.3pt}\end{center}

\emph{Astronomisches Jahrbuch für das Jahr} 1817 \emph{nebst einer Sammlung der neuesten in die
astronomischen Wissenschaften einschlagenden Abhandlungen, Beobachtungen und Nachrichten.
Berechnet und herausgegeben von} J. E. \textsc{Bode}, \emph{Königl. Astronomen und Mitglied der
Academie.} Beim Verfasser und in Commission bei J. E. \textsc{Hitzig}. 
Berlin 1814. 260 Seiten in Octav, nebst einer Kupfertafel.

Das Jahr 1817 hat nur zwei in Europa nicht sichtbare Sonnenfinsternisse. Unter den Phänomenen,
die die besondere Aufmerksamkeit der Astronomen verdienen, zeichnen wir eine sehr nahe
Zusammenkunft der Venus mit dem Regulus aus, die nach Hrn. \textsc{Bode}'s Berechnung den
28. September $14^{\mathrm{u}}$ Berliner Zeit Statt haben wird, und sich demnach in dem
östlichen Theile von Europa sehr gut wird beobachten lassen. Nach Hrn. \textsc{Bode}'s Rechnung
geht die Venus nur Eine Minute dem Stern südlich vorbei, so dass zu einer wirklichen Bedeckung
nur sehr wenig fehlen wird. Dergleichen Phänomene verdienen im Voraus nach aller Schärfe berechnet
zu werden, und \emph{wirkliche} Bedeckungen der grössern Sterne von der Venus, besonders wenn
diese näher bei der untern Conjunction mit der Sonne sich befindet, verdienten unsrer Ueberzeugung
nach eben so sehr, dass ihretwegen grosse Reisen unternommen würden, wie die Durchgänge der Venus
vor der Sonne. Merkwürdig ist auch eine eben so nahe aber in Europa unsichtbare Zusammenkunft des
Mars mit dem Saturn den 18. April $7^{\mathrm{u}}$, die der Sternwarte auf Botanybay zur
Beobachtung zu empfehlen sein würde.

Die \emph{Zusätze} zu dem vorliegenden Bande des Jahrbuches fangen mit Ephemeriden für die Pallas
von Hrn. \textsc{Nicolai}, für die Juno von Hrn. \textsc{Möbius} und für die Vesta von Hrn.
\textsc{Gerling} an. Es folgen Beobachtungen und parabolische Elemente des zweiten Cometen vom
Jahre 1813 von Hrn. Dr. \textsc{Olbers}. --- Astronomische Beobachtungen auf der Prager Sternwarte
im Jahre 1813 von den Hrn. \textsc{David} und \textsc{Bittner}. Die hier empfohlene Methode,
Azimuthe durch correspondirende horizontale Abstände von einem Sterne bei gleichen Höhen ohne
Zuziehung der Zeit zu bestimmen, ist an sich sehr gut, und würde auch, bei Anwendung eines
Instruments, an dem beide, der horizontale und der verticale Kreis hinlängliche Grösse und
Eintheilung haben, zu sehr genauen Bestimmungen geeignet sein. Allein bei Anwendung eines
Theodolithen, dessen Höhenkreis, oder eines Kreises, dessen Azimuthalkreis nur in Minuten
getheilt ist, kann natürlich nur eine beschränktere Genauigkeit davon erwartet werden. Bei dem
erstern Instrument würde jenem nachtheiligen Umstande zwar ausgewichen, wenn man das Fernrohr auf
der beobachteten Höhe unverrückt stehen liesse, dann fiele aber dagegen der Vortheil der
Vervielfältigung weg. Bei Anwendung des Kreises würden ausserdem noch kleine Fehler in der
Berichtigung der Querlibelle und des Parallelismus der Gesichtslinie mit der Ebene des
Instruments, (die nicht wohl ganz zu vermeiden, aber bei der eigentlichen Bestimmung des Instruments
unschädlich sind) in ihrer ganzen Stärke auf das Resultat wirken: diese Fehler lassen sich
indessen aufheben, wenn man am folgenden Tage dieselben Beobachtungen bei einer entgegengesetzten
Lage des Kreises wiederholt, da sie sodann
\SourceLeaf{582}auch in entgegengesetztem Sinn wirken müssen. --- Noch einige Bemerkungen und Zusätze zu den Beobachtungen über $\eta$ Antinous im astronomischen Jahrbuche 1816 von Hrn. Prof. \textsc{Wurm}. Mit Zuziehung von zehn ältern Beobachtungen von \textsc{Pigott}, welche Hrn. \textsc{Wurm} bisher entgangen waren, und noch einigen eigenen, so wie mit Anwendung der Methode der kleinsten Quadrate bestimmte dieser geschickte Astronom die Periode nunmehro auf 7,17604 Tage, und die Epoche 1800 Januar 4,504. Zugleich werden nach dieser Bestimmung neue Tafeln mitgetheilt. --- Verschiedene schätzbare astronomische Beiträge von Hrn. Prof. \textsc{Littrow} in Kasan betreffen die Entwickelung der Mittelpunktsgleichung in eine Reihe; die Theorie der Epicykel, durch die, wie der Verf. sehr gut zeigt, die elliptischen Bewegungen zugleich in Beziehung auf die Länge und den Radius Vector sich nicht darstellen lassen; eine Methode die Kreisbahn eines Planeten aus zwei Beobachtungen zu bestimmen, welche zwar allgemeiner und der Form nach einfacher ist, als die von demselben Verfasser im Jahrbuch 1816 entwickelte; aber bei mässigen Bewegungen wegen des eingeführten Cosinus derselben in der Ausübung geringere Genauigkeit geben würde, obwohl freilich gegenwärtig die ganze Aufgabe nur ein theoretisches Interesse hat; ferner (wovon dieselbe Bemerkung gilt) Bestimmung einer geradlinigen Bahn aus vier beobachteten Längen (dieser Aufsatz ist nur durch Druckfehler ganz entstellt); dieselbe Aufgabe für drei vollständige Beobachtungen, wo der Verf. die aus der \emph{Theoria motus corporum coelestium} entlehnte Auflösung auf drei Vestabeobachtungen anwendet (an sich ist die Voraussetzung wie in dem angeführten Werke, Art. 131 gezeigt ist, auch nicht einmal als Näherung statthaft, wiewohl Herr \textsc{Littrow} sie auch nur vorschlägt, um einen neu entdeckten Himmelskörper nach kurzer Unterbrechung leichter wieder zu finden, zu welchem Zweck indessen auch schon die gewöhnlichen Interpolationsmethoden hinreichen werden). Ferner giebt Herr \textsc{Littrow} eine indirecte Auflösung der Aufgabe, aus den Höhen zweier Sterne die Zeit und Polhöhe zu finden, die zwar an sich zweckmässig ist, aber an Bequemlichkeit noch gewinnt, wenn man sich dabei der logarithmischen Differenzen bedient; ausserdem muss auch noch bemerkt werden, dass so oft die Stundenwinkel $\alpha$, $\alpha'$ kleiner als $90^{\circ}$ sind, man besser thut, $\frac{1}{2}\alpha$ und $\frac{1}{2}\alpha'$ durch ihre Sinus nach den bekannten Formeln zu bestimmen. Die neue Methode, Circummeridianhöhen zu berechnen, würden wir der gewöhnlichen nachsetzen; die Vergrösserung der Genauigkeit ist zu unbedeutend, da doch der Fehler von derselben Ordnung bleibt, wie bei der letztern, wenn man sich mit dem ersten Gliede begnügt. Eben so ziehen wir bei der letzten auf das Erdsphäroid sich beziehenden Aufgabe den hier gegebenen Reihen die ihnen gleichgültigen endlichen Ausdrücke vor. --- Es folgen hierauf astronomische Beobachtungen auf der Sternwarte zu Wien von den Hrn. \textsc{Triesnecker} und \textsc{Bürg} und auf der Sternwarte zu Cremsmünster von Hrn. \textsc{Derfflinger}. --- Beobachtungen zur Bestimmung der geographischen Lage von Port Jackson, aus \textsc{Malaspinas} Papieren berechnet von Hrn. J. \textsc{Oltmanns}. Es wird uns hier Hoffnung gemacht zu einer vollständigen Bearbeitung sämmtlicher Beobachtungen. --- Die Ideen zur Perturbationsrechnung nach \textsc{Keppler} von Hrn. Prof. J. W. \textsc{Pfaff} in Nürnberg, scheinen uns doch zu oberflächlich, um irgend ein Resultat zu geben. Die hier aufgestellte Form für die Störungsgleichungen des Knoten und der Neigung (oder vielmehr für ihre Differentialänderungen) scheint uns ganz unfruchtbar, und es folgen namentlich keine solche Beziehungen daraus, wie Herr \textsc{Pfaff} glaubt, da z. B. Gleichungen von einerlei Periode aus den Gliedern $Fxx'$ und $Gyy'$ hervorgehen. --- Beobachtungen auf der Wilnaer Sternwarte von Hrn. Prof. \textsc{Sniadecky}. --- Bemerkungen über angestellte geographische Ortsbestimmungen in Ungarn, Oesterreich und Bayern, vom Hrn. Prof. \textsc{Bürg}, enthalten eine Rechtfertigung gegen einen Aufsatz der monatlichen Correspondenz, in welchem man besonders die zu Wien gemachten Azimuthalbestimmung zweifelhaft zu machen gesucht hatte. --- Ueber den Cometen 1558 von Hrn. Dr. \textsc{Olbers}. Man kannte von diesem Cometen bisher nur drei Beobachtungen des Landgrafen \textsc{Wilhelm} vom 20., 21. und 23. September, und eine des \textsc{Cornelius Gemma} vom 20. September, welche letztere von
\SourceLeaf{583}der gleichzeitigen des Landgrafen sehr verschieden war. Durch eine sehr glückliche Verbesserung des
Textes bei \textsc{Gemma}, indem statt distabat 20 die gelesen wird, distabat eo die hebt Herr Dr. \textsc{Olbers} diesen Widerspruch auf das befriedigendste. \textsc{Gemma}'s Beobachtung gilt nun für den 17. August, und macht
eine wenigstens einigermassen genäherte Bestimmung der Bahn des Cometen, welche hier mitgetheilt
wird, möglich. --- Einige physisch-astronomische Beobachtungen des Saturn, Mars, des Mondes, der
Venus und Sonne von Hrn. Dr. \textsc{Gruithusen} in München geben merkwürdige Beweise von der Vortrefflichkeit der \textsc{Fraunhofer}'schen Fernröhre. --- Noch von Hrn. Prof. \textsc{Oltmanns} ein Beitrag zur Längenbestimmung von \textsc{Quito} aus verschiedenen von den Französischen Academikern gemachten Beobachtungen;
und die Berechnung der unlängst zur Sprache gebrachten totalen Sonnenfinsterniss zu Mirabeau in der
Provence den 3. Juni 1239. Herr \textsc{Oltmanns} macht davon eine glückliche Anwendung auf die Bewegung
des Mondsknoten, die so wie sie in Hrn. \textsc{Oltmanns} Tafeln angesetzt war, hienach eine Verringerung
von $1'40''$ auf 100 Jahre nöthig hat. Mit Zuziehung von zwei andern ähnlichen Beobachtungen findet
Herr \textsc{Oltmanns} im Mittel diese Verminderung der hundertjährigen Bewegung $1'14''$. --- Die Tafeln für
die scheinbaren Oerter des Polarsterns von Hrn. Prof. \textsc{Bessel}, sind für die beobachtenden Astronomen
ein ungemein schätzbares Geschenk. Es ist schade, dass der Verfasser dabei die neuern Beobachtungen
der geraden Aufsteigung dieses Sterns von verschiedenen Astronomen noch nicht benutzen konnte, die
übereinstimmend die Nothwendigkeit einer Vergrösserung (von etwa 4 Secunden in Zeit gegenwärtig)
beweisen. Wir hätten gewünscht, dass auch die reinen mittlern Stellungen für jedes Jahr beigefügt
wären, um die Bestimmungen anderer Astronomen bequemer vergleichen zu können. Von Hrn. Prof.
\textsc{Bürg} Bemerkungen über die Revision seiner frühern Mondberechnungen. Hr. \textsc{Bürg} hat angefangen, zum
Behuf einer nochmaligen Verbesserung seiner Mondtafeln, die Greenwicher Beobachtungen seit 1765
mit denselben zu vergleichen: die Resultate dieser mühsamen Arbeit, wie sie auch immer ausfallen
mögen, werden gewiss interessant und lehrreich sein. --- Beobachtungen der Juno und neue Elemente
ihrer Bahn; Berechnung der nächsten Opposition der Pallas und andere astronomische Nachrichten von
Hrn. Prof. \textsc{Gauss}. --- Etwas über die Erwartung neuer Entdeckungen am Himmel durch Fernröhre. Es
werden hier die Schwierigkeiten entwickelt, die der fortschreitenden Verbesserung der Sehwerkzeuge
Grenzen setzen: wir gestehen indess, dass wir dadurch doch nicht von der völligen Unüberwindlichkeit
dieser Hindernisse überzeugt sind. In wie fern die Brauchbarkeit der Vorschläge, das Fernrohr durch
ein Uhrwerk der Bewegung der Sterne folgen zu lassen, durch die hier hervorgehobenen Mängel gemindert oder gar aufgehoben werde, müssen wir auf sich beruhen lassen, auch verstehen wir nicht ganz,
wie S.\,218 behauptet werden kann, dass bei dem beständigen Hin- und Herfliegen des Bildes eine ruhige Wahrnehmung \emph{während der wenigen Secunden seines Durchganges} unmöglich sei, da ja eben die
Vorrichtung die Wirkung hat, das Bild auf lange Zeit im Felde zu erhalten. Sollte indess die Erfahrung das Gewicht dieses Vorwurfs bestätigen, so würde, däucht uns, die Mechanik wohl Mittel finden
können, auch ohne Zahn und Getriebe. dem Fernrohr eine höchst sanfte und gleichförmige Bewegung
zu geben. Das Haupthinderniss, die Wirkung der Sehwerkzeuge immer weiter zu treiben, möchte wohl
in den atmosphärischen Oscillationen liegen, die desto merklicher werden, je mehr jene vergrössern. ---
Ueber zwei veränderliche Sterne im Herkules von Hrn. Dr. \textsc{Koch} in Danzig. Leider ist auch die Danziger Sternwarte, aus der so manche nützliche Beobachtungen hervorgegangen sind, ein Opfer der letzten Belagerung geworden. Der in der Anmerkung von Hrn. \textsc{Bode} mit $m$ bezeichnete Stern, der noch
in keinem Verzeichnisse steht, ist übrigens schon von unserm Hrn. Prof. \textsc{Harding} bestimmt worden, und
befindet sich auf dem XVI. Blatt von dessen Himmelskarten. --- Die mittlern Stellungen von 28 der
vornehmsten Sterne aus den Plejaden nach \textsc{Piazzi}'s neuem Sternkatalog; aus eben demselben die mittlern
Stellungen der 36 \textsc{Maskelyne}'schen Fundamentalsterne, und endlich die des Polarsterns. Sehr merkwür\SourceLeaf{584}dig ist bei letzterm die aus den Beobachtungen des Hrn \textsc{Cacciatore} folgende Rectascensionsparallaxe von
$2^{\mathrm{s}}885$ in Zeit, die uns indessen zu gross scheint, und sich mit den zahlreichen von Hrn. \textsc{von Lindenau}
auf der Seeberger Sternwarte angestellten Beobachtungen nicht vereinigen lässt. Aus \textsc{Oriani}'s beobachteten
Zenithdistanzen hatten wir nur eine halb so grosse Parallaxe abgeleitet (M. s. das 10. St. dieser Blätter
[S. 574 d. B.]), die freilich auch keine absolute Zuverlässigkeit hat. Die jährliche eigne Bewegung des
Polarsterns in gerader Aufsteigung findet \textsc{Piazzi} aus Vergleichung

\begin{center}
\small\begin{tabular}{@{}l@{\hspace{8pt}}c@{\hspace{8pt}}r@{}}
mit Hevel & $.\;.\;.\;.\;.\;.\;.\;.\;.\;.\;.\;.\;.$ & $+6^{\mathrm{s}}82$\\
Flamstead & $.\;.\;.\;.\;.\;.\;.\;.\;.\;.\;.\;.$ & $+9.03$\\
La Caille & $.\;.\;.\;.\;.\;.\;.\;.\;.\;.\;.\;.$ & $+3.96$\\
Bradley & $.\;.\;.\;.\;.\;.\;.\;.\;.\;.\;.\;.$ & $+1.62$\\
\end{tabular}
\end{center}
\noindent in Zeit; Herr \textsc{von Zach} hatte aus Vergleichung der Bestimmungen von \textsc{La Caille} und \textsc{Mathieu} gefunden
$+3^{\mathrm{s}}177$. Es wäre sehr zu wünschen, dass dieser wichtige Gegenstand einmal sorgfältig untersucht
würde, wobei aber alle Originalbeobachtungen vorher auf eine gleichförmige Art reducirt werden müssten. ---
Astronomische Beobachtungen auf der Königl Sternwarte zu Berlin im Jahre 1813. --- Beobachtungen auf der
Sternwarte zu Dorpat von Herrn Prof. \textsc{Struve} in den Jahren 1812, 1813 und 1814. --- Noch verschiedene
Beobachtungen zu Cremsmünster. --- Ferner Nachrichten über den Doppelstern 61 Cygni aus \textsc{Piazzi}'s Libro
sesto del reale osservatorio (wo die merkwürdige eigne Bewegung dieses Doppelsterns zuerst angezeigt war)
und aus dem neuen grossen Sterncatalog dieses Astronomen. --- Nachweisung, dass von acht am Himmel vermissten
Fixsternen keiner die Ceres, Pallas, Juno oder Vesta war. Wir bemerken hiebei, dass alle diese Sterne bis auf
den in der nordlichen Krone, bereits im neunten Bande der monatlichen Correspondenz (S. 153. 151. 241. 155.
243. 241. 246) discutirt, und namentlich alle hier angeführten \textsc{Mayer}'schen Sterne aus den Originalpapieren
bereits befriedigend erklärt sind. ---

Unter den vermischten Bemerkungen und Nachrichten, die wie gewöhnlich am Schlusse des Jahrbuches angehängt
sind, heben wir hier nur eine aus, nemlich eine Stelle aus einem Briefe des Hrn. Dr. \textsc{Reche} in Mühlheim
am Rhein vom 6. Mai 1814, welcher zufolge Herr Prof. \textsc{Kramp} in Strassburg während der Blockade dieser
Stadt das Problem, die Bahn eines Cometen aus drei Beobachtungen zu bestimmen, auf eine neue Weise aufgelöst
und gefunden hat, dass alle von ihm berechneten Cometenbahnen, die man bisher für Parabeln angesehen hatte,
\emph{ganz gewiss und bestimmt} Hyperbeln sind. Da eine solche Ankündigung sehr auffallend sein muss, und die
Arbeit des Hrn. \textsc{Kramp} seitdem im Julistück der mathematischen Zeitschrift erschienen ist, die zu Nimes
unter dem Titel Annales de mathématiques pures et appliquées seit einigen Jahren herauskommt, so wird es
hier der Ort sein, mit der vorstehenden Anzeige noch die Beurtheilung dieser Abhandlung zu verbinden.

Es kommt hier auf zwei Umstände an, theils auf den Werth der Methode an sich, theils auf die Art ihrer
Anwendung auf wirklich beobachtete Cometen, und die Zuverlässigkeit der daraus hervorgehenden Resultate.
Wir wollen zuerst einiges über die letzteren bemerken. In dieser Abhandlung ist nur von der Anwendung auf
Einen Cometen die Rede, nemlich auf den von 1781, wo Herr \textsc{Kramp} aus den Beobachtungen vom 14, 19. und
25. November eine Hyperbel ableitet, deren Excentricität $= 4,580652$, und Periheldistanz $1,048364$. Wer
nur einige Erfahrung in Berechnung von Cometenbahnen hat, dem würde die kurze Zwischenzeit schon hinreichen,
um die Zuverlässigkeit dieses Resultats zu würdigen. Wir müssen indessen noch hinzusetzen, dass weit
entfernt, die gefundenen Elemente an andern Beobachtungen zu prüfen, sie nicht einmal mit den drei zum
Grunde liegenden Beobachtungen selbst verglichen werden, und dass eine solche Prüfung demjenigen, der sie
etwa selbst vornehmen wollte, dadurch wenigstens sehr erschwert ist, dass Herr \textsc{Kramp} die Berechnung
der Elemente nicht einmal zu
\SourceLeaf{585}Ende geführt hat, so dass Länge des Perihel und Durchgangszeit fehlen. Wenn man nun endlich noch bemerkt,
dass die Erdbahn bei der ganzen Rechnung als ein Kreis betrachtet ist, so wird man es wohl für eine
überflüssige Mühe halten, noch besonders zu untersuchen, ob vielleicht auch in der numerischen
Rechnung selbst noch Fehler begangen sind, wofür der \emph{bedeutende} Unterschied zwischen dem Werthe der
Grösse
\[
P\!\cdot\!Q'' - P''Q = 0.4614411
\]
und der Summe der beiden Grössen
\[
\begin{aligned}
PQ' - P'Q &= 0.2133547\\
P'Q'' - P''Q' &= 0.2495918
\end{aligned}
\]
welche Summe eigentlich, als Folge der Methode, jener ersten Grössen \emph{genau} gleich werden sollte, zu
sprechen scheint. (Herr \textsc{Kramp} urtheilt hierüber anders, und nennt diese beträchtliche Verschiedenheit
eine egalité presque rigoureuse, aber wozu mit sieben Decimalen rechnen, wenn man nicht einmal die dritte
verbürgen könnte?) Wir sagen \emph{scheint}, weil es vielleicht möglich wäre, dass dieser Unterschied
eine Folge der Weglassungen S. 13 wäre.

Wir wenden uns nun zu der Methode selbst. Es ist schon erinnert, dass Herr \textsc{Kramp} (in der That ohne
alle Noth) die Excentricität der Erdbahn ganz vernachlässigt. Das Wesentliche der Methode lässt sich
mit wenigen Worten angeben. Herr \textsc{Kramp} bezeichnet die halbe grosse Axe der Bahn des Himmelskörpers
mit $b$; die Coordinaten seiner drei Oerter in der Bahn und auf die Knotenlinie als Abscissenlinie bezogen
mit $bP$, $bP'$, $bP''$; $bQ$, $bQ'$, $bQ''$. Diese sechs Grössen lassen sich durch die Beobachtungsdata,
die Neigung der Bahn $\beta$ und die Länge des aufsteigenden Knoten $\delta$ (wie bekannt ist) leicht
bestimmen. Von der andern Seite findet hingegen Herr \textsc{Kramp}, durch Combination der Formeln die sich
auf die \textsc{Kepler}'schen Gesetze beziehen, und indem er die Unterschiede der excentrischen Anomalien
ihren Sinus gleich setzt, die drei Grössen
\[
\begin{gathered}
PQ' - P'Q\\
P'Q'' - P''Q'\\
PQ'' - P''Q
\end{gathered}
\]
resp. den Zwischenzeiten $t$, $t'$, $t+t'$ proportional, so dass man die Gleichung hat
\[
0 = PQ' - P'Q + P'Q'' - P''Q' + P''Q - PQ''
\]
welche Herr \textsc{Kramp} eine équation essentielle et très remarquable par sa simplicité nennt. ---
Indem nun Herr \textsc{Kramp} in den beiden Gleichungen
\[
\begin{aligned}
(t+t')(PQ'-P'Q) &= t(PQ''-P''Q)\\
(t+t')(P'Q''-P''Q') &= t'(PQ''-P''Q)
\end{aligned}
\]
für $P$, $P'$, $P''$, $Q$, $Q'$, $Q''$ ihre Werthe durch $\beta$ und $\delta$ substituirt (wo $b$ offenbar
herausfällt), erhält er zwei Gleichungen zwischen $\beta$ und $\delta$, aus denen er mit Weglassung von
ein Paar Gliedern, durch Elimination von $\beta$ endlich eine kubische Gleichung für $\operatorname{tang}\delta$ ableitet.
Aus dieser wird $\delta$ bestimmt, sodann aus einer der beiden eben gedachten Gleichungen (oder beiden)
$\beta$, sodann die Grössen $bP$, $bP'$, $bP''$, $bQ$, $bQ'$, $bQ''$, und aus vier dieser Grössen, nach
Methoden die wir Kürze halber hier übergehen, die übrigen Elemente des Kegelschnittes.

Es wird sehr leicht sein, die gänzliche Unstatthaftigkeit dieser Methode zu beweisen. Wir sehen hier an
die Stelle der Bewegungsgesetze die Proportionalität der mehrmals erwähnten drei Grössen zu den
Zwischenzeiten treten. Blos auf diese Proportionalität oder die ihr gleichgeltenden zwei Gleichun\SourceLeaf{586}gen und den Satz, dass die drei Oerter in Einer Ebene liegen, wird die ganze Auflösung gegründet. Die
drei Grössen
\[
bb(PQ'-P'Q), \quad bb(P'Q''-P''Q'), \quad bb(PQ''-P''Q)
\]
sind, wie man sogleich bemerken wird, nichts anderes als die drei in der Theoria motus Corporum Coelestium
art. 112. mit $n''$, $n$, $n'$ bezeichneten Grössen. Es ist aber dort im Art. 114. bewiesen, dass die
Distanzen von der Erde \emph{linearisch} bestimmt werden, sobald das Verhältniss dieser drei Grössen als
bekannt angesehen wird. Es ist folglich ein Fehler, wenn Herr \textsc{Kramp} eine Aufgabe, deren \emph{strenge}
Auflösung blos eine Gleichung des ersten Grades erfordert, auf eine cubische Gleichung bringt, selbst mit
Vernachlässigung einiger Glieder, mit deren Zuziehung bei Hrn. \textsc{Kramp}'s Behandlung sogar eine
biquadratische Gleichung hervorgegangen wäre. Doch dies ist noch weniger wichtig, als der zweite Umstand,
dass in dem angeführten Werke im Art. 131. schon ausführlich die Unstatthaftigkeit bewiesen ist, jene
drei Grössen $n''$, $n$, $n'$ \emph{zugleich} den Zwischenzeiten $t$, $t'$, $t+t'$ proportional zu setzen,
indem das Resultat auch nicht einmal als Annäherung gelten darf.

War Hrn. \textsc{Kramp}, wie es fast scheint, die Theoria motus corporum coelestium, welche fünf Jahre vor
der Blokade von Strassburg erschienen ist, unbekannt, so hätte er sich doch leicht überzeugen können,
dass der wahre Geist seiner Methode eigentlich auf einem Princip beruhet, dessen Verwerflichkeit schon seit
langer Zeit bekannt ist. Die oben erwähnten drei Grössen sind nemlich, wie einige Aufmerksamkeit sogleich
zeigt, nichts weiter als die doppelten Flächen der Dreiecke zwischen dem ersten und zweiten, zweiten und
dritten, ersten und dritten Radius Vector; diese drei Grössen also \emph{zugleich} den drei Zwischenzeiten
proportional setzen, heisst nichts weiter, als die drei Oerter des Cometen in Einer geraden Linie annehmen,
deren beide Stücke den Zwischenzeiten proportional sind. Und so ist auch Hrn. \textsc{Kramp}'s sogenannte
équation essentielle durchaus nichts weiter, als die Bedingungsgleichung für die Lage der drei Oerter in
Einer geraden Linie. Dass aber in dem Cometenproblem diese Voraussetzung theils auf eine Auflösung führen
muss, die nur Gleichungen des ersten Grades erfordert, theils aber wirklich auch nicht einmal als Näherung
zulässig ist, dies sind längst bekannte Wahrheiten, worüber man auch besonders in \textsc{Olbers}
vortrefflicher und von Hrn. \textsc{Kramp} selbst in gegenwärtiger Abhandlung angeführten Schrift einen
eben so gründlichen als lichtvollen Beweis findet.
